% ---------------------------------------------------------------------------------------
% MODELO DE TCC - INSTITUTO FEDERAL DE SERGIPE (IFS)
% ---------------------------------------------------------------------------------------
% Autor do Modelo: Mikael dos Santos Gois (mikaelgois1@gmail.com)
% Baseado nos modelos de: 
%   - Diogo D. Moreira (TCC, Dissertação e Tese)
%   - DECOM/UFOP (Monografia de Ciência da Computação)
% Modelo Base: abnTeX2 adaptado para as normativas internas de TCC do IFS (Art. 31 e 34)
% Versão: 1.0 (2026)
% ---------------------------------------------------------------------------------------
% Este arquivo faz parte de um modelo opensource para estudantes do IFS.
% Não remova os créditos dos autores originais e do modelo base.
% Este modelo é fornecido "no estado em que se encontra", sem garantias de qualquer tipo.
% ---------------------------------------------------------------------------------------

% Lista de símbolos:
% Elemento opcional. Deve ser elaborada de acordo com a ordem apresentada no texto, com o devido 
% significado. - NBR 14724:2024 (4.2.1.12)

\begin{simbolos}
  \item[$\alpha$] Coeficiente de atrito
  \item[$\beta$] Ângulo de inclinação
  \item[$\gamma$] Razão de compressão
  \item[$\delta$] Deslocamento
  \item[$\epsilon$] Erro relativo
  \item[$\zeta$] Fator de amortecimento
  \item[$\eta$] Eficiência
  \item[$\theta$] Temperatura
  \item[$\lambda$] Comprimento de onda
  \item[$\mu$] Coeficiente de viscosidade
  \item[$\pi$] Número pi (3.14159...)
  \item[$\rho$] Densidade
  \item[$\sigma$] Desvio padrão
  \item[$\tau$] Tempo de relaxação
  \item[$\phi$] Potencial elétrico
  \item[$\chi$] Susceptibilidade magnética
  \item[$\psi$] Função de onda
  \item[$\omega$] Frequência angular
\end{simbolos}